\section{Data Provenance and Supplementary Protocols}
\label{appendix:supplementary}
\label{app:provenance}

This appendix supports the reproducibility of the \emph{review itself}: it documents the sources of the reported figures consolidated in Sections~\ref{sec:reported}--\ref{sec:deployment}, and provides the reference prompt templates and quantization-pipeline specifications under which those figures were reported. (The \gls{qlora} adaptation recipe appears in Section~\ref{sec:reported}, Table~\ref{tab:hyperparam_search}; detailed device specifications appear in Section~\ref{sec:deployment}, Table~\ref{tab:appendix_hardware}.) The appendix does not report original experimental runs.

\subsection{Provenance of the Consolidated Figures}
\label{subsec:appendix_provenance_table}

% -----------------------------------------------------------------------------
% AUTHOR ACTION REQUIRED BEFORE SUBMISSION.
% Table \ref{tab:provenance_models} below records model-level provenance, which is
% complete. Per-cell attribution for Tables \ref{tab:finetuning},
% \ref{tab:edge_deployment}, \ref{tab:energy} and \ref{tab:latency_compare} is NOT
% yet reproduced anywhere in this manuscript. Either (a) supply the per-cell
% source ledger for those four tables, or (b) withdraw them and the claims that
% depend on them. Section \ref{subsec:appendix_provenance_status} states the
% current position honestly, but it is not a substitute for (a) or (b).
% -----------------------------------------------------------------------------

Provenance is recorded at two levels. \emph{Model-level} provenance fixes, for
each model in the cohort, the primary release document from which its
architectural specification and headline benchmark scores derive
(Table~\ref{tab:provenance_models}). \emph{Source-class} provenance records, for
each consolidated table, the categories of source the figures are drawn from and
the reconciliation rules applied
(Section~\ref{subsec:appendix_provenance_status}). Benchmark definitions and
metric variants follow the originating task papers cited in
Section~\ref{subsec:reported_benchmarks}.

\begin{table}[ht]
\centering
\caption{Model-level provenance for the reviewed cohort: the primary release document from which each model's architectural specification and headline reported scores derive.}
\label{tab:provenance_models}
\footnotesize\adjustbox{max width=\linewidth}{%
\begin{tabular}{llll}
\toprule
\textbf{Model(s)} & \textbf{Developer} & \textbf{Primary source} & \textbf{Source type} \\
\midrule
TinyLlama-1.1B & --- & \citet{zhang2024tinyllama} & Technical report \\
SmolLM2-1.7B, -360M & Hugging Face & \citet{allal2025smollm2} & Technical report \\
StableLM 2-1.6B & Stability AI & \citet{stabilityai2024stablelm} & Technical report \\
Qwen2.5-0.5B/1.5B/3B/7B & Alibaba & \citet{qwen2024qwen25} & Technical report \\
Llama 3.2-1B/3B & Meta & \citet{grattafiori2024llama3} & Technical report \\
Gemma 2-2B/9B & Google DeepMind & \citet{gemma2team2024} & Technical report \\
Gemma 3-1B/4B & Google DeepMind & \citet{team2025gemma3} & Technical report \\
Mistral 7B & Mistral AI & \citet{jiang2023mistral} & Technical report \\
Phi-3-mini & Microsoft & \citet{abdin2024phi3} & Technical report \\
Phi-3.5-mini & Microsoft & \citet{abdin2024phi35} & Technical report \\
Phi-4-mini & Microsoft & \citet{microsoft2025phi4mini} & Technical report \\
Phi-4 & Microsoft & \citet{abdin2024phi4} & Technical report \\
Ministral 8B & Mistral AI & Vendor model card & Model card \\
\bottomrule
\end{tabular}%
}
\end{table}

\subsection{Provenance Status by Table}
\label{subsec:appendix_provenance_status}

The consolidated tables do not all rest on an equally firm reporting base, and
the distinction matters for how much weight each should carry.

\emph{Architectural and accuracy figures.} The specifications in
Table~\ref{tab:model_specs} and the accuracy figures in
Tables~\ref{tab:zero_shot} and~\ref{tab:appendix_zero_shot_extended} derive from
the primary release documents in Table~\ref{tab:provenance_models}, supplemented
where a model's report omits a task by standardized community evaluations
(the Open \gls{llm} Leaderboard and \texttt{lm-evaluation-harness}-based reports
\citep{gao2023harness}). These are the best-supported figures in the review.

\emph{Compression and quantization figures.} The degradation figures in
Table~\ref{tab:quantization} and the footprints in
Table~\ref{tab:memory_footprint} derive from the quantization-method literature
\citep{frantar2023gptq,lin2024awq,kim2024squeezellm} and from published
quantized artifact sizes. The comparative characteristics in
Table~\ref{tab:compression_compare} are a synthesis of the compression
literature reviewed in Section~\ref{sec:compression} rather than measurements of
a single cohort, and the retention column should be read as an indicative band
rather than a measured quantity.

\emph{Adaptation, efficiency, and energy figures.} Tables~\ref{tab:finetuning},
\ref{tab:edge_deployment}, \ref{tab:energy} and~\ref{tab:latency_compare} rest on
a materially thinner base. Public reporting of fine-tuning gains, edge
throughput, and per-inference energy is sparse, is rarely produced under a common
protocol, and in several cases exists only for a subset of the cohort; some
entries are consequently aggregates or interpolations across sources rather than
single quoted values, and this manuscript does not reproduce a per-cell ledger
for them. They are offered as indicative orders of magnitude that establish
relative position, not as measurements, and no claim in this paper should be
read as resting on their precise values. Converting them into measured
quantities is the first item on the verification agenda of
Section~\ref{subsec:challenges_verification}.

The remaining subsections of this appendix document the reference conditions---prompt formats, quantization toolchains, adaptation hyperparameters---that the reviewed sources most commonly report; these are preserved here so that readers can interpret any individual reported figure against the conditions it was measured under.

\subsection{Prompt Formats}
\label{subsec:appendix_prompt_templates}
The reviewed source materials commonly preserve each model's own instruction
wrapper while following a shared task-payload structure, which is summarized here
in prose. \textbf{Multiple-choice question answering} presents the question
followed by the enumerated options (A--D) and an ``Answer:'' cue, with the model
expected to emit a single option letter. \textbf{Mathematical reasoning} prepends
an instruction to solve the problem step by step and appends a ``Let's think step
by step'' cue, so that a chain-of-thought precedes the final numeric answer.
\textbf{Extractive question answering} supplies the context passage, then the
question, then an instruction to answer from the context and to respond
``unanswerable'' when the context does not contain the answer. \textbf{Abstractive
summarization} supplies the source article preceded by an instruction to produce a
two- to three-sentence summary. Where a source reports few-shot results, the same
fixed set of exemplars is prepended to the payload for every model on a given
task. Because sources differ in the exact wording of these wrappers and in
answer-extraction heuristics, the consolidated figures are compared only within
tasks and, wherever possible, within a single reporting source.

\subsection{Quantization and Compression Pipeline}
\label{subsec:appendix_quantization_pipeline}
To complement the methodology section, the compression workflows reported across the reviewed sources are summarized here. The reviewed literature reports five precision or compression settings across the consolidated comparison:
\begin{description}
    \item[FP16 baseline] Standard half-precision inference used as the reference condition.
    \item[GGUF Q4\_K\_M] Mixed-precision 4-bit quantization via \texttt{llama.cpp}, assigning higher precision to the most sensitive tensors.
    \item[AWQ 4-bit] Activation-aware weight quantization with group size 128 and calibration on held-out text sequences.
    \item[GPTQ 4-bit] Layer-wise post-training quantization using second-order curvature approximations.
    \item[bitsandbytes NF4] 4-bit NormalFloat quantization with double quantization, matching the \gls{qlora} training regime.
\end{description}

For compression stress tests, the sources report SparseGPT~\citep{frantar2023sparsegpt}-induced 50\% unstructured sparsity applied before 4-bit quantization. Reported figures for this combined pruning-plus-quantization setting are retained in the provenance tables for sensitivity analysis rather than the main comparative tables, because support for sparse inference remains hardware dependent and is not uniformly reported across the reviewed sources.

\subsection{Edge Measurement Protocols Reported in the Source Literature}
\label{subsec:appendix_edge_protocol}
The deployment figures consolidated in Sections~\ref{sec:reported}--\ref{sec:deployment} are drawn from source materials whose stated measurement protocols typically include a multi-minute warm-up to stabilize thermal and frequency behavior before measurement. Across the reviewed edge-deployment sources, the reported per-model, per-device measurements cover:
\begin{description}
    \item[Time-to-first-token (\gls{ttft})] Reported latency from prompt submission to first generated token, typically averaged over 50 prompts with 128-token inputs.
    \item[Sustained throughput] Reported mean tokens-per-second over 256-token generations.
    \item[Peak memory] Reported maximum RAM or unified-memory consumption during generation.
    \item[Load time] Reported elapsed time from process start until the model is ready to serve requests.
    \item[Energy per inference] Reported device-level power estimates where the source provides counter- or telemetry-based measurements.
\end{description}
Because these reporting conventions differ in small but potentially consequential ways across sources (warm-up duration, prompt-length distribution, power-counter granularity), readers should treat the consolidated figures as indicative rather than precisely comparable; the conflict-resolution rules in Section~\ref{subsec:reported_sources} are applied whenever discrepant protocols yield materially different reported values for the same tuple.

\subsection{Extended Results Tables}
\label{subsec:appendix_extended_results}
The representative five-shot results for the four models with fine-tuning
comparisons, and the QLoRA hyperparameter configuration used to produce them, are
given with the fine-tuning comparison itself in Section~\ref{sec:reported}
(Table~\ref{tab:finetuning}, Table~\ref{tab:hyperparam_search}) rather than
repeated here.

\begin{landscape}
\begin{table}[ht]
\centering
\caption{Extended zero-shot results across a broader task set (FP16, instruction-tuned variants). All values are reported figures.}
\label{tab:appendix_zero_shot_extended}
\footnotesize\adjustbox{max width=\linewidth}{%
\begin{tabular}{lcccccccccccc}
\toprule
\textbf{Model} & \textbf{ARC-E} & \textbf{ARC-C} & \textbf{HellaSwag} & \textbf{TriviaQA} & \textbf{SQuAD 2.0} & \textbf{GSM8K} & \textbf{MATH} & \textbf{HumanEval} & \textbf{MBPP} & \textbf{MMLU} & \textbf{SST-2} & \textbf{MNLI} \\
\midrule
Phi-4-mini-Inst. & 92.1 & 83.7 & 69.1 & 68.5 & 74.2 & 88.6 & 64.0 & 72.0 & 74.8 & 67.3 & 92.4 & 72.1 \\
Qwen2.5-7B-Inst. & 94.8 & 90.1 & 80.0 & 72.4 & 76.8 & 88.7 & 60.4 & 84.8 & 79.2 & 72.6 & 93.8 & 74.6 \\
Gemma 2-9B-Inst. & 95.2 & 89.8 & 87.1 & 71.8 & 75.5 & 84.9 & 51.3 & 68.9 & 73.5 & 71.3 & 93.5 & 73.8 \\
Gemma 3-4B-Inst. & 92.4 & 85.2 & 78.4 & 64.2 & 70.8 & 82.1 & 53.7 & 64.6 & 69.3 & 65.8 & 91.2 & 70.4 \\
LLaMA 3.2-3B-Inst. & 88.5 & 76.1 & 77.2 & 58.6 & 65.4 & 75.6 & 46.7 & 58.5 & 62.4 & 61.8 & 88.2 & 63.5 \\
LLaMA 3.2-1B-Inst. & 78.2 & 62.3 & 63.8 & 42.5 & 52.8 & 47.2 & 24.1 & 35.4 & 40.2 & 48.6 & 82.5 & 54.2 \\
Mistral 7B-Inst. v0.3 & 91.8 & 82.7 & 73.5 & 65.8 & 68.4 & 81.9 & 41.6 & 66.5 & 68.3 & 63.0 & 90.8 & 68.2 \\
Ministral 8B-Inst. & 93.2 & 85.4 & 76.8 & 67.2 & 71.2 & 83.5 & 45.2 & 69.5 & 71.8 & 65.7 & 92.1 & 71.5 \\
Qwen2.5-3B-Inst. & 90.2 & 78.6 & 72.1 & 58.8 & 66.5 & 79.8 & 50.2 & 65.2 & 67.8 & 65.6 & 89.5 & 66.8 \\
Qwen2.5-0.5B-Inst. & 68.4 & 52.1 & 48.3 & 28.5 & 38.2 & 38.5 & 18.7 & 22.0 & 28.4 & 47.5 & 78.2 & 48.5 \\
SmolLM2-1.7B-Inst. & 74.5 & 59.8 & 60.2 & 35.4 & 52.3 & 42.6 & 19.3 & 28.7 & 35.6 & 50.1 & 82.6 & 54.2 \\
SmolLM2-360M-Inst. & 55.8 & 38.2 & 42.5 & 18.2 & 32.4 & 15.8 & 5.2 & 8.5 & 12.3 & 33.4 & 72.4 & 42.8 \\
TinyLlama 1.1B & 62.4 & 47.5 & 53.8 & 28.8 & 42.5 & 18.4 & 6.8 & 12.2 & 16.8 & 36.2 & 76.5 & 46.2 \\
StableLM 2 1.6B & 72.8 & 55.2 & 57.6 & 32.5 & 48.6 & 32.1 & 12.4 & 22.5 & 28.9 & 44.3 & 80.2 & 52.5 \\
\bottomrule
\end{tabular}%
}
\end{table}
\end{landscape}

\subsection{Human-Evaluation Conventions in the Source Literature}
\label{subsec:appendix_human_eval}
For subjective tasks such as summarization quality and instruction following, the reviewed sources that report human judgments typically rate outputs along three axes---\textbf{relevance}, \textbf{coherence}, and \textbf{factual accuracy}---on a 1--5 Likert scale, with model identity hidden from annotators and inter-annotator reliability reported via a Krippendorff-style coefficient. MT-Bench-style results substitute a strong \gls{llm} judge for human raters. Because these protocols differ across sources in rubric wording and judge model---and because per-model coverage is sparse---instruction-following and summarization are not tabulated in this review and are discussed only qualitatively.
